\section{Involutive object ledger}
\label{sec:involutive_ledger}

The data of the paper are a measured space of objects with an
involution, and an observable compatible with it.

\begin{definition}[Involutive object ledger]
\label{def:duality_confinement:involutive-object-ledger}
An \emph{involutive object ledger} is a tuple
\[
  \IOL=(X,J,\mu,\psi,Y,\Jiso)
\]
consisting of
\begin{enumerate}
  \item a measurable space \(X\), the \emph{carrier};
  \item an involution \(J:X\to X\), that is, \(J(Jx)=x\) for every
    \(x\in X\);
  \item a positive measure \(\mu\) on \(X\);
  \item a real Hilbert space \(Y\), the \emph{response space};
  \item a linear isometric involution \(\Jiso:Y\to Y\), that is,
    \(\Jiso^2=I_Y\) and
    \(\langle\Jiso v,\Jiso w\rangle=\langle v,w\rangle\) for all
    \(v,w\in Y\);
  \item a strongly measurable map \(\psi:X\to Y\), the
    \emph{readout}, that is \emph{equivariant}:
    \[
      \psi(Jx)=\Jiso\psi(x)
      \qquad (x\in X).
    \]
\end{enumerate}
The \emph{fixed locus} of the ledger is
\(\Fix(J)=\{x\in X: Jx=x\}\).
\end{definition}

The carrier-side involution \(J\) acts on the objects whose mass is
being studied, while \(\Jiso\) acts on the space in which those objects
are observed; equivariance says that observing \(Jx\) is the same as
observing \(x\) and then applying \(\Jiso\).  Strong measurability
(being a pointwise limit of measurable simple functions) is the
standard condition under which Bochner integrals of \(Y\)-valued maps
are defined; when \(Y\) is separable it is equivalent to Borel
measurability.  Every result below remains true if \(\psi\) is only
\(\mu\)-almost everywhere strongly measurable (measurability statements
then hold up to \(\mu\)-null sets), and the formal verification uses
that weaker hypothesis.  The measure \(\mu\) is not assumed
\(\sigma\)-finite; under the standing hypothesis
\(\int_X\|\psim\|^2\,d\mu<\infty\) of the later sections, the set
\(\{\psim\neq0\}\) is automatically \(\sigma\)-finite.  No measurability of \(J\) and no invariance of \(\mu\)
under \(J\) is assumed: none of the results below needs them.

\paragraph{The finite weighted case.}
If \(X\) is a finite set with all subsets measurable and every point
has finite mass, the measure is a family of real weights
\(\mu(x)\geq0\), and integrals are finite sums; this is the
\emph{finite weighted case}.  The \emph{visible support} of \(\mu\) is the set of objects of
positive weight.  In general, a statement holds for \(\mu\)-almost
every \(x\) when the set where it fails is contained in a measurable
set of measure zero; in the finite case this means that it holds at
every visible object.  The mass of the ledger is \emph{confined to the
fixed locus} when \(Jx=x\) for \(\mu\)-almost every \(x\), that is,
when \(X\setminus\Fix(J)\) is \(\mu\)-null.  If \(\Fix(J)\) is
measurable this says \(\mu(X\setminus\Fix(J))=0\).

\begin{example}[Complex conjugation]
Let \(X\subset\mathbb{C}\) be a Borel set stable under complex
conjugation, with \(\mu\) any positive Borel measure on \(X\).  Put
\[
  J(z)=\overline{z},\qquad
  Y=\mathbb{R}^2,\qquad
  \Jiso(u,v)=(u,-v),\qquad
  \psi(z)=(\operatorname{Re}z,\operatorname{Im}z).
\]
Then \(\psi(Jz)=(\operatorname{Re}z,-\operatorname{Im}z)=\Jiso\psi(z)\),
so the readout is equivariant.  The fixed locus is \(X\cap\mathbb{R}\),
and the coordinate that changes sign is the imaginary part.  Mass is
confined to the fixed locus exactly when \(\mu\) is carried by the real
axis.
\end{example}

\begin{example}[Zeros of the Riemann zeta function]
The same pattern underlies the motivating example.  The map
\(J_L(s)=1-\overline{s}\) is an involution of \(\mathbb{C}\) with fixed
locus the critical line \(\operatorname{Re}s=\tfrac12\).  Taking for
\(X\) the nontrivial zeros of the Riemann zeta function, \(\mu\) the
counting measure weighted by multiplicity (the zeros are symmetric
under \(J_L\) by the functional equation and the reflection principle
\citep{TitchmarshHeathBrown1986}), \(Y=\mathbb{R}\),
\(\Jiso=-I\) and \(\psi(s)=\operatorname{Re}s-\tfrac12\), one obtains an
involutive object ledger whose mass is confined to the fixed locus
exactly when the Riemann hypothesis holds.  Here \(\psi\) is already
odd, so \(\psim=\psi\).  This example only illustrates the definition.
The results below also require
\(\int_X\|\psim\|^2\,d\mu=\sum_\rho(\operatorname{Re}\rho-\tfrac12)^2<\infty\),
summed over the zeros with multiplicity; the sum vanishes if the
Riemann hypothesis holds, but its finiteness is not known otherwise.
Nothing in this paper bears on the location of the zeros.
\end{example}
